\documentclass[11pt,letterpaper]{report}
\usepackage[margin=1in]{geometry}
\usepackage{booktabs}
\usepackage{xltabular}
\usepackage[expansion=false]{microtype}
\usepackage{parskip}
\usepackage[colorlinks=true, linkcolor=black, urlcolor=blue]{hyperref}
\renewcommand{\arraystretch}{1.25}
\usepackage{titlesec}
\titleformat{\chapter}[display]
  {\normalfont\Huge\bfseries}{}{0pt}{\Huge}
\titlespacing*{\chapter}{0pt}{10pt}{20pt}
\title{\textbf{SVX Library Class Catalog}\\[0.5em]
       \large SystemVerilog Extension Library --- Core Library (\texttt{src/})\\[0.3em]
       \normalsize Version 2.1.0 beta}
\author{Mark Glasser}
\date{September 2026}
\begin{document}
\maketitle
\tableofcontents

\chapter{Types}
\section{svx\_types --- Standard Typedefs}
\noindent\small\textit{File: \texttt{src/types/svx\_types.svh}}\par

Standardized C-style integer typedefs, size and index types, process identifiers, and package-level Boolean constants. Provides a portable numeric vocabulary across the library.

\begin{xltabular}{\linewidth}{>{\small}p{1.5in} >{\small}p{1.5in} >{\small}X}
\toprule
\textbf{Typedef} & \textbf{Base Type} & \textbf{Description} \\
\midrule
\endhead
\bottomrule
\endfoot
\texttt{int8\_t} & \texttt{byte} & Signed 8-bit integer \\
\texttt{uint8\_t} & \texttt{byte unsigned} & Unsigned 8-bit integer \\
\texttt{int16\_t} & \texttt{shortint} & Signed 16-bit integer \\
\texttt{uint16\_t} & \texttt{shortint unsigned} & Unsigned 16-bit integer \\
\texttt{int32\_t} & \texttt{int} & Signed 32-bit integer \\
\texttt{uint32\_t} & \texttt{int unsigned} & Unsigned 32-bit integer \\
\texttt{int64\_t} & \texttt{longint} & Signed 64-bit integer \\
\texttt{uint64\_t} & \texttt{longint unsigned} & Unsigned 64-bit integer \\
\texttt{int128\_t} & \texttt{bit [127:0]} & Signed 128-bit integer (two-state) \\
\texttt{uint128\_t} & \texttt{bit unsigned [127:0]} & Unsigned 128-bit integer (two-state) \\
\texttt{size\_t} & \texttt{uint64\_t} & Container/array size \\
\texttt{index\_t} & \texttt{uint64\_t} & Unsigned container index \\
\texttt{signed\_index\_t} & \texttt{int64\_t} & Signed container index (for skip/offset) \\
\texttt{pid\_t} & \texttt{int unsigned} & Process identifier \\
\end{xltabular}

\section{type\_handle}
\noindent\small\textit{File: \texttt{src/types/type\_handle.svh}}\par

Singleton type-handle mechanism. \texttt{type\_handle\#(T)} is a parameterized class that returns a unique singleton object for each distinct type \texttt{T}, enabling runtime type identity comparison. \texttt{type\_container\#(T)} wraps any value with its type handle for polymorphic storage.

\subsection*{\texttt{type\_handle\_base}}
\noindent\small\textit{virtual class}\quad \textbf{extends} \texttt{object}\par

Non-parameterized base for all type handles.

\begin{xltabular}{\linewidth}{>{\small}p{1.35in} >{\small}p{1.25in} >{\small}p{0.95in} >{\small}X}
\toprule
\textbf{Method} & \textbf{Return} & \textbf{Args} & \textbf{Description} \\
\midrule
\endhead
\bottomrule
\endfoot
\texttt{get\_type\_handle} & \texttt{type\_handle\_base} & --- & Return the singleton type handle for this object (pure virtual). \\
\end{xltabular}

\subsection*{\texttt{type\_handle\#(T)}}
\noindent\small\textit{class}\quad \textbf{extends} \texttt{type\_handle\_base}\par

Parameterized singleton. One instance exists per distinct type T.

\begin{xltabular}{\linewidth}{>{\small}p{1.35in} >{\small}p{1.25in} >{\small}p{0.95in} >{\small}X}
\toprule
\textbf{Method} & \textbf{Return} & \textbf{Args} & \textbf{Description} \\
\midrule
\endhead
\bottomrule
\endfoot
\texttt{get\_type} & \texttt{type\_handle\#()} & --- & Return (or lazily create) the singleton for type T. Static. \\
\texttt{get\_type\_handle} & \texttt{type\_handle\_base} & --- & Return this singleton cast to base type. \\
\end{xltabular}

\subsection*{\texttt{type\_container\_base}}
\noindent\small\textit{class}\par

Non-parameterized base for typed containers; supports polymorphic storage with convert2string.

\begin{xltabular}{\linewidth}{>{\small}p{1.35in} >{\small}p{1.25in} >{\small}p{0.95in} >{\small}X}
\toprule
\textbf{Method} & \textbf{Return} & \textbf{Args} & \textbf{Description} \\
\midrule
\endhead
\bottomrule
\endfoot
\texttt{convert2string} & \texttt{string} & --- & Return a human-readable representation (default: empty string). \\
\texttt{get\_type\_handle} & \texttt{type\_handle\_base} & --- & Return the type handle for the contained type (default: null). \\
\end{xltabular}

\subsection*{\texttt{type\_container\#(T)}}
\noindent\small\textit{class}\quad \textbf{extends} \texttt{type\_container\_base}\par

Holds a value of type T along with its type handle.

\noindent\textbf{Constructor:} \texttt{new(---)} --- Initialises the internal type handle singleton for T.

\begin{xltabular}{\linewidth}{>{\small}p{1.35in} >{\small}p{1.25in} >{\small}p{0.95in} >{\small}X}
\toprule
\textbf{Method} & \textbf{Return} & \textbf{Args} & \textbf{Description} \\
\midrule
\endhead
\bottomrule
\endfoot
\texttt{get\_type\_handle} & \texttt{type\_handle\_base} & --- & Return the type handle for T. \\
\texttt{set} & void & \texttt{T x} & Store a new value, replacing any existing one. \\
\texttt{get} & \texttt{T} & --- & Retrieve the stored value. \\
\end{xltabular}

\section{type\_match}
\noindent\small\textit{File: \texttt{src/types/type\_match.svh}}\par

Compile-time type-compatibility utilities. \texttt{type\_match\#(T1,T2)} provides static functions to check whether two type parameters are identical or whether one is derived from the other. The convenience macro \texttt{\`{}check\_is\_derived\_from} installs a const-initializer that will call \texttt{\$fatal} if the constraint is violated at elaboration time.

\subsection*{\texttt{type\_match\#(T1, T2)}}
\noindent\small\textit{class}\par

Static type-comparison utilities for T1 and T2.

\begin{xltabular}{\linewidth}{>{\small}p{1.35in} >{\small}p{1.25in} >{\small}p{0.95in} >{\small}X}
\toprule
\textbf{Method} & \textbf{Return} & \textbf{Args} & \textbf{Description} \\
\midrule
\endhead
\bottomrule
\endfoot
\texttt{test\_is\_match} & \texttt{bit} & --- & True if T1 and T2 are the same type. Static. \\
\texttt{is\_match} & \texttt{bit} & \texttt{int line, string file} & True if match; fatal error with location if not. Static. \\
\texttt{test\_is\_derived\_from} & \texttt{bit} & --- & True if T1 is derived from T2 (via \$cast probe). Static. \\
\texttt{is\_derived\_from} & \texttt{bit} & \texttt{int line, string file} & True if derived; fatal error if not. Static. \\
\texttt{fail\_match} & void & \texttt{int line, string file} & Issue \$fatal reporting type mismatch. Static. \\
\texttt{fail\_derived} & void & \texttt{int line, string file} & Issue \$fatal reporting derivation failure. Static. \\
\end{xltabular}

\section{typeid}
\noindent\small\textit{File: \texttt{src/types/typeid.svh}}\par

Runtime type-classification for common SystemVerilog scalar types. \texttt{typeid\#(T)} uses \texttt{type\_handle} singletons to determine at run time whether \texttt{T} is integral, two-state, four-state, real, or string. Classification is limited to explicitly enumerated built-in types.

\subsection*{\texttt{typeid\#(T)}}
\noindent\small\textit{class}\par

Static type-classification functions for built-in SV scalar types.

\begin{xltabular}{\linewidth}{>{\small}p{1.35in} >{\small}p{1.25in} >{\small}p{0.95in} >{\small}X}
\toprule
\textbf{Method} & \textbf{Return} & \textbf{Args} & \textbf{Description} \\
\midrule
\endhead
\bottomrule
\endfoot
\texttt{is\_int} & \texttt{bit} & --- & True if T is any integral type. Static. \\
\texttt{is\_int\_fail} & \texttt{bit} & --- & Same as is\_int(); prints error and calls \$finish if false. Static. \\
\texttt{is\_two\_state} & \texttt{bit} & --- & True if T is integral and two-state (e.g., bit, int). Static. \\
\texttt{is\_two\_state\_fail} & \texttt{bit} & --- & Same but terminates on failure. Static. \\
\texttt{is\_four\_state} & \texttt{bit} & --- & True if T is logic, reg, integer, or time. Static. \\
\texttt{is\_four\_state\_fail} & \texttt{bit} & --- & Same but terminates on failure. Static. \\
\texttt{is\_real} & \texttt{bit} & --- & True if T is real. Static. \\
\texttt{is\_real\_fail} & \texttt{bit} & --- & Same but terminates on failure. Static. \\
\texttt{is\_string} & \texttt{bit} & --- & True if T is string. Static. \\
\texttt{is\_string\_fail} & \texttt{bit} & --- & Same but terminates on failure. Static. \\
\end{xltabular}

\chapter{Traits}
\section{Traits Classes}
\noindent\small\textit{File: \texttt{src/containers/traits.svh}}\par

The traits system provides per-type constants and static functions used by all parameterized containers and algorithms. Every traits class defines an \texttt{empty\_t} typedef, an \texttt{empty} constant, and static \texttt{equal()}, \texttt{compare()}, and \texttt{sort()} functions. The hierarchy root is \texttt{void\_traits}; scalar integer families derive from \texttt{base\_int\_traits}.

\subsection*{\texttt{void\_traits}}
\noindent\small\textit{class}\quad \textbf{extends} \texttt{void\_t}\par

Root traits class. empty is null (void\_t). All type-flag localparams default to false.

\begin{xltabular}{\linewidth}{>{\small}p{1.35in} >{\small}p{1.25in} >{\small}p{0.95in} >{\small}X}
\toprule
\textbf{Method} & \textbf{Return} & \textbf{Args} & \textbf{Description} \\
\midrule
\endhead
\bottomrule
\endfoot
\texttt{equal} & \texttt{bit} & \texttt{void\_t a, b} & Always returns 1 (void objects are equivalent). Static. \\
\texttt{compare} & \texttt{int32\_t} & \texttt{void\_t a, b} & Returns 0 if equal (always), else 1. Static. \\
\texttt{sort} & void & \texttt{ref void\_t vec[\$]} & No-op sort. Static. \\
\texttt{check\_trait} & \texttt{bit} & \texttt{bit trait, bit exp, int line, string file} & Assert trait==exp; call fail() if not. Static. \\
\texttt{fail} & void & \texttt{int line, string file} & Issue \$fatal with optional file/line. Static. \\
\end{xltabular}

\subsection*{\texttt{object\_traits}}
\noindent\small\textit{class}\quad \textbf{extends} \texttt{void\_traits}\par

Traits for object-derived classes. Delegates equal()/compare() to object::compare().

\begin{xltabular}{\linewidth}{>{\small}p{1.35in} >{\small}p{1.25in} >{\small}p{0.95in} >{\small}X}
\toprule
\textbf{Method} & \textbf{Return} & \textbf{Args} & \textbf{Description} \\
\midrule
\endhead
\bottomrule
\endfoot
\texttt{equal} & \texttt{bit} & \texttt{object a, b} & True if a.compare(b)==0. Static. \\
\texttt{compare} & \texttt{int32\_t} & \texttt{object a, b} & Returns a.compare(b). Static. \\
\texttt{sort} & void & \texttt{ref object vec[\$]} & No-op (objects are unordered). Static. \\
\end{xltabular}

\subsection*{\texttt{class\_traits\#(T)}}
\noindent\small\textit{class}\quad \textbf{extends} \texttt{void\_traits}\par

Traits for arbitrary class handles. empty is null. Uses identity (==) for equality.

\begin{xltabular}{\linewidth}{>{\small}p{1.35in} >{\small}p{1.25in} >{\small}p{0.95in} >{\small}X}
\toprule
\textbf{Method} & \textbf{Return} & \textbf{Args} & \textbf{Description} \\
\midrule
\endhead
\bottomrule
\endfoot
\texttt{equal} & \texttt{bit} & \texttt{T a, T b} & True if a==b (handle identity). Static. \\
\texttt{compare} & \texttt{int32\_t} & \texttt{T a, T b} & Returns 0 if equal, else 1. Static. \\
\texttt{sort} & void & \texttt{ref T vec[\$]} & No-op. Static. \\
\end{xltabular}

\subsection*{\texttt{base\_int\_traits\#(T)}}
\noindent\small\textit{class}\quad \textbf{extends} \texttt{void\_traits}\par

Parameterized base traits for integral scalar types. empty=T'(0). Provides arithmetic compare and built-in sort.

\begin{xltabular}{\linewidth}{>{\small}p{1.35in} >{\small}p{1.25in} >{\small}p{0.95in} >{\small}X}
\toprule
\textbf{Method} & \textbf{Return} & \textbf{Args} & \textbf{Description} \\
\midrule
\endhead
\bottomrule
\endfoot
\texttt{equal} & \texttt{bit} & \texttt{T a, T b} & True if a==b. Static. \\
\texttt{compare} & \texttt{int32\_t} & \texttt{T a, T b} & Returns 1 if a>b, -1 if a<b, 0 if equal. Static. \\
\texttt{sort} & void & \texttt{ref T vec[\$]} & Sort using SystemVerilog built-in queue sort. Static. \\
\end{xltabular}

\subsection*{\texttt{base\_unsigned\_int\_traits\#(T)}}
\noindent\small\textit{class}\quad \textbf{extends} \texttt{base\_int\_traits\#(T)}\par

Sets is\_signed=0, is\_unsigned=1. Convenience base for all unsigned integer traits typedefs.


\subsection*{\texttt{base\_signed\_int\_traits\#(T)}}
\noindent\small\textit{class}\quad \textbf{extends} \texttt{base\_int\_traits\#(T)}\par

Sets is\_signed=1, is\_unsigned=0. Convenience base for all signed integer traits typedefs.


\subsection*{\texttt{base\_four\_state\_traits\#(T)}}
\noindent\small\textit{class}\quad \textbf{extends} \texttt{base\_int\_traits\#(T)}\par

Traits base for four-state types (logic, reg, integer, time). Sets is\_two\_state=0, is\_four\_state=1.


\subsection*{\texttt{integer\_traits}}
\noindent\small\textit{class}\quad \textbf{extends} \texttt{base\_four\_state\_traits\#(integer)}\par

Traits for the SystemVerilog integer type (four-state, signed).


\subsection*{\texttt{reg\_traits / logic\_traits / time\_traits}}
\noindent\small\textit{class}\quad \textbf{extends} \texttt{base\_four\_state\_traits\#(reg)}\par

Traits for reg, logic, and time (all four-state, unsigned). logic\_traits and time\_traits are typedefs of reg\_traits.


\subsection*{\texttt{bit\_vector\_traits\#(N)}}
\noindent\small\textit{class}\quad \textbf{extends} \texttt{base\_int\_traits\#(bit[N-1:0])}\par

Traits for packed bit-vector types of width N. is\_array=1; is\_arithmetic=1.


\subsection*{\texttt{real\_base\_traits\#(T)}}
\noindent\small\textit{class}\quad \textbf{extends} \texttt{void\_traits}\par

Traits for real/shortreal. Uses epsilon-based equality (2.2e-16). Sorted with item comparator.

\begin{xltabular}{\linewidth}{>{\small}p{1.35in} >{\small}p{1.25in} >{\small}p{0.95in} >{\small}X}
\toprule
\textbf{Method} & \textbf{Return} & \textbf{Args} & \textbf{Description} \\
\midrule
\endhead
\bottomrule
\endfoot
\texttt{equal} & \texttt{bit} & \texttt{T a, T b} & True if |a-b| <= epsilon. Static. \\
\texttt{compare} & \texttt{int32\_t} & \texttt{T a, T b} & 1 if a>b (and not equal), -1 if a<b, else 0. Static. \\
\texttt{sort} & void & \texttt{ref T vec[\$]} & Sort with item comparator. Static. \\
\end{xltabular}

\subsection*{\texttt{string\_traits}}
\noindent\small\textit{class}\quad \textbf{extends} \texttt{void\_traits}\par

Traits for string. empty="". Lexicographic compare and sort.

\begin{xltabular}{\linewidth}{>{\small}p{1.35in} >{\small}p{1.25in} >{\small}p{0.95in} >{\small}X}
\toprule
\textbf{Method} & \textbf{Return} & \textbf{Args} & \textbf{Description} \\
\midrule
\endhead
\bottomrule
\endfoot
\texttt{equal} & \texttt{bit} & \texttt{string a, b} & True if a==b. Static. \\
\texttt{compare} & \texttt{int32\_t} & \texttt{string a, b} & Lexicographic: 1 if a>b, -1 if a<b, 0 if equal. Static. \\
\texttt{sort} & void & \texttt{ref string vec[\$]} & Sort with item comparator. Static. \\
\end{xltabular}

\chapter{Containers}
\section{object --- Base Object}
\noindent\small\textit{File: \texttt{src/containers/object.svh}}\par

Defines \texttt{void\_t} (the empty base class, used as the "void" type throughout SVX) and \texttt{object} (adds virtual to\_str, compare, and copy to void\_t). All containers and tree nodes derive ultimately from \texttt{object}.

\subsection*{\texttt{void\_t}}
\noindent\small\textit{class}\par

The empty class. Serves as the null base type and as the type of the "empty" element for void\_traits.


\subsection*{\texttt{object}}
\noindent\small\textit{class}\quad \textbf{extends} \texttt{void\_t}\par

Base object providing virtual to\_str(), compare(), and copy().

\begin{xltabular}{\linewidth}{>{\small}p{1.35in} >{\small}p{1.25in} >{\small}p{0.95in} >{\small}X}
\toprule
\textbf{Method} & \textbf{Return} & \textbf{Args} & \textbf{Description} \\
\midrule
\endhead
\bottomrule
\endfoot
\texttt{to\_str} & \texttt{string} & --- & Return a human-readable string (default: ""). \\
\texttt{compare} & \texttt{int} & \texttt{object obj} & Compare to obj; return 0 (default implementation). \\
\texttt{copy} & void & \texttt{object rhs} & Shallow-copy from rhs (default: no-op). \\
\end{xltabular}

\section{container --- Container Base Classes}
\noindent\small\textit{File: \texttt{src/containers/container.svh}}\par

Abstract container base hierarchy. \texttt{container} defines size/clear/is\_empty. \texttt{typed\_container\#(T,P)} adds the protected \texttt{m\_empty} member initialised from \texttt{P::empty} at construction.

\subsection*{\texttt{container}}
\noindent\small\textit{virtual class}\quad \textbf{extends} \texttt{void\_t}\par

Abstract base for all containers. Defines the core size/clear/is\_empty interface.

\begin{xltabular}{\linewidth}{>{\small}p{1.35in} >{\small}p{1.25in} >{\small}p{0.95in} >{\small}X}
\toprule
\textbf{Method} & \textbf{Return} & \textbf{Args} & \textbf{Description} \\
\midrule
\endhead
\bottomrule
\endfoot
\texttt{size} & \texttt{size\_t} & --- & Return the number of elements (pure virtual). \\
\texttt{clear} & void & --- & Remove all elements (pure virtual). \\
\texttt{is\_empty} & \texttt{bit} & --- & True if container holds no elements (pure virtual). \\
\end{xltabular}

\subsection*{\texttt{typed\_container\#(T, P)}}
\noindent\small\textit{virtual class}\quad \textbf{extends} \texttt{container}\par

Typed abstract base. Initialises m\_empty from P::empty in the constructor.

\noindent\textbf{Constructor:} \texttt{new(---)} --- Casts P::empty to T and stores in protected m\_empty.

\begin{xltabular}{\linewidth}{>{\small}p{1.35in} >{\small}p{1.25in} >{\small}p{0.95in} >{\small}X}
\toprule
\textbf{Method} & \textbf{Return} & \textbf{Args} & \textbf{Description} \\
\midrule
\endhead
\bottomrule
\endfoot
\texttt{size} & \texttt{size\_t} & --- & Overridable; returns 0 by default. \\
\texttt{is\_empty} & \texttt{bit} & --- & Overridable; returns 0 by default. \\
\texttt{clear} & void & --- & Overridable; no-op by default. \\
\end{xltabular}

\section{vector}
\noindent\small\textit{File: \texttt{src/containers/vector.svh}}\par

Dynamic array backed by a SystemVerilog queue. Provides random read/write with automatic extension, shallow copy/clone, equality comparison, sort, and bulk append operations. The base class for queue, stack, and deque.

\subsection*{\texttt{vector\#(T, P)}}
\noindent\small\textit{class}\quad \textbf{extends} \texttt{typed\_container\#(T,P)}\par

Ordered, random-access dynamic array.

\begin{xltabular}{\linewidth}{>{\small}p{1.35in} >{\small}p{1.25in} >{\small}p{0.95in} >{\small}X}
\toprule
\textbf{Method} & \textbf{Return} & \textbf{Args} & \textbf{Description} \\
\midrule
\endhead
\bottomrule
\endfoot
\texttt{create} & \texttt{this\_t} & \texttt{vector\_type list} & Static factory: create a vector pre-loaded from a queue literal. \\
\texttt{size} & \texttt{size\_t} & --- & Number of elements currently in the vector. \\
\texttt{is\_empty} & \texttt{bit} & --- & True when size()==0. \\
\texttt{clear} & void & --- & Delete all elements; reset to empty state. \\
\texttt{extend} & void & \texttt{size\_t sz} & Push sz empty elements, growing the vector by that amount. \\
\texttt{write} & void & \texttt{index\_t idx, T t} & Write element at idx; auto-extends if idx >= size(). \\
\texttt{read} & \texttt{T} & \texttt{index\_t idx} & Read element at idx; returns m\_empty if out of range. \\
\texttt{copy} & void & \texttt{this\_t vec} & Shallow copy from vec into this vector. \\
\texttt{clone} & \texttt{this\_t} & --- & Return a new vector that is a shallow copy of this one. \\
\texttt{compare} & \texttt{int} & \texttt{this\_t v} & Returns 0 if equal to v, else non-zero. \\
\texttt{equal} & \texttt{bit} & \texttt{this\_t v} & Element-wise equality using P::equal(). \\
\texttt{sort} & void & --- & Sort elements using P::sort(). \\
\texttt{appendv} & void & \texttt{vector\_type v} & Append a queue literal to this vector. \\
\texttt{appendc} & void & \texttt{T t} & Append a single element to this vector. \\
\texttt{append} & void & \texttt{this\_t v} & Append all elements of vector v to this vector. \\
\end{xltabular}

\section{map}
\noindent\small\textit{File: \texttt{src/containers/map.svh}}\par

Class-based dynamic associative array. Allows sparse key-value storage allocated on demand and passed by reference. \texttt{singleton\_map} provides a global instance accessible via \texttt{get\_inst()}.

\subsection*{\texttt{map\#(KEY, T, P)}}
\noindent\small\textit{class}\quad \textbf{extends} \texttt{typed\_container\#(T,P)}\par

Associative array keyed by KEY, storing values of type T.

\begin{xltabular}{\linewidth}{>{\small}p{1.35in} >{\small}p{1.25in} >{\small}p{0.95in} >{\small}X}
\toprule
\textbf{Method} & \textbf{Return} & \textbf{Args} & \textbf{Description} \\
\midrule
\endhead
\bottomrule
\endfoot
\texttt{get} & \texttt{T} & \texttt{KEY key} & Return item for key; returns m\_empty if key not found. \\
\texttt{insert} & \texttt{bit} & \texttt{KEY key, T item} & Insert (key,item). Returns 1 if key is new, 0 if duplicate (value overwritten). \\
\texttt{size} & \texttt{size\_t} & --- & Number of key-value pairs. \\
\texttt{is\_empty} & \texttt{bit} & --- & True when size()==0. \\
\texttt{delete} & \texttt{bit} & \texttt{KEY key} & Remove entry for key. Returns 1 on success, 0 if not found. \\
\texttt{clear} & void & --- & Remove all entries. \\
\texttt{exists} & \texttt{bit} & \texttt{KEY key} & True if key exists in the map. \\
\texttt{copy} & void & \texttt{this\_t rhs} & Shallow copy from rhs into this map. \\
\texttt{clone} & \texttt{this\_t} & --- & Return a new map that is a shallow copy of this one. \\
\texttt{compare} & \texttt{int32\_t} & \texttt{this\_t m} & Returns 0 if maps are equal, else non-zero. \\
\texttt{equal} & \texttt{bit} & \texttt{this\_t m} & True if all keys and P::equal()-checked values match. \\
\texttt{first} & \texttt{bit} & \texttt{ref KEY index} & Move to first key (smallest); returns 0 if empty. For iterator use. \\
\texttt{last} & \texttt{bit} & \texttt{ref KEY index} & Move to last key (largest). For iterator use. \\
\texttt{next} & \texttt{bit} & \texttt{ref KEY index} & Advance to next key. For iterator use. \\
\texttt{prev} & \texttt{bit} & \texttt{ref KEY index} & Move to previous key. For iterator use. \\
\end{xltabular}

\subsection*{\texttt{singleton\_map\#(KEY, T, P)}}
\noindent\small\textit{class}\quad \textbf{extends} \texttt{map\#(KEY,T,P)}\par

A map of which only one instance exists per type parameterization.

\begin{xltabular}{\linewidth}{>{\small}p{1.35in} >{\small}p{1.25in} >{\small}p{0.95in} >{\small}X}
\toprule
\textbf{Method} & \textbf{Return} & \textbf{Args} & \textbf{Description} \\
\midrule
\endhead
\bottomrule
\endfoot
\texttt{get\_inst} & \texttt{this\_t} & --- & Return the singleton instance (lazily created). Static. \\
\end{xltabular}

\section{queue}
\noindent\small\textit{File: \texttt{src/containers/queue.svh}}\par

FIFO queue derived from vector. put() inserts at the head; get() removes from the tail. \texttt{fixed\_size\_queue} adds a maximum-size constraint and a push-success status.

\subsection*{\texttt{queue\#(T, P)}}
\noindent\small\textit{class}\quad \textbf{extends} \texttt{vector\#(T,P)}\par

FIFO queue: items enter at the front and leave from the back.

\begin{xltabular}{\linewidth}{>{\small}p{1.35in} >{\small}p{1.25in} >{\small}p{0.95in} >{\small}X}
\toprule
\textbf{Method} & \textbf{Return} & \textbf{Args} & \textbf{Description} \\
\midrule
\endhead
\bottomrule
\endfoot
\texttt{put} & void & \texttt{T t} & Push item onto the head (front) of the queue. \\
\texttt{get} & \texttt{T} & --- & Pop and return item from the tail; returns empty if queue is empty. \\
\texttt{peek} & \texttt{T} & --- & Return the tail item without removing it. \\
\texttt{is\_empty} & \texttt{bit} & --- & True when size()==0. \\
\texttt{clone} & \texttt{this\_t} & --- & Return a shallow copy of this queue. \\
\end{xltabular}

\subsection*{\texttt{fixed\_size\_queue\#(T, P)}}
\noindent\small\textit{class}\quad \textbf{extends} \texttt{queue\#(T,P)}\par

Queue that rejects pushes when the maximum size is reached.

\noindent\textbf{Constructor:} \texttt{new(size\_t n=1)} --- Set maximum capacity to n (minimum 1).

\begin{xltabular}{\linewidth}{>{\small}p{1.35in} >{\small}p{1.25in} >{\small}p{0.95in} >{\small}X}
\toprule
\textbf{Method} & \textbf{Return} & \textbf{Args} & \textbf{Description} \\
\midrule
\endhead
\bottomrule
\endfoot
\texttt{set\_max\_size} & void & \texttt{size\_t n} & Update the capacity limit; clamps to current size if n is too small. \\
\texttt{get\_max\_size} & \texttt{size\_t} & --- & Return the current capacity limit. \\
\texttt{put} & void & \texttt{T t} & Push item; silently rejected if at capacity. \\
\texttt{last\_push\_succeeded} & \texttt{bit} & --- & True if the most recent put() succeeded. \\
\texttt{is\_full} & \texttt{bit} & --- & True when size() >= max\_size. \\
\texttt{copy} & void & \texttt{this\_t q} & Copy contents and state (max\_size, push\_ok) from q. \\
\texttt{clone} & \texttt{this\_t} & --- & Return a shallow copy. \\
\texttt{equal} & \texttt{bit} & \texttt{vector\#() v} & True if element contents, max\_size, and push\_ok all match. \\
\end{xltabular}

\section{stack}
\noindent\small\textit{File: \texttt{src/containers/stack.svh}}\par

LIFO stack derived from vector. push() and pop() operate on the front (index 0).

\subsection*{\texttt{stack\#(T, P)}}
\noindent\small\textit{class}\quad \textbf{extends} \texttt{vector\#(T,P)}\par

LIFO stack: top of stack is index 0.

\begin{xltabular}{\linewidth}{>{\small}p{1.35in} >{\small}p{1.25in} >{\small}p{0.95in} >{\small}X}
\toprule
\textbf{Method} & \textbf{Return} & \textbf{Args} & \textbf{Description} \\
\midrule
\endhead
\bottomrule
\endfoot
\texttt{push} & void & \texttt{T t} & Push item onto the top of the stack. \\
\texttt{pop} & \texttt{T} & --- & Remove and return the top item; returns empty if stack is empty. \\
\texttt{peek} & \texttt{T} & --- & Return top item without removing it. \\
\texttt{is\_empty} & \texttt{bit} & --- & True when size()==0. \\
\texttt{clone} & \texttt{this\_t} & --- & Return a shallow copy of this stack. \\
\end{xltabular}

\section{deque}
\noindent\small\textit{File: \texttt{src/containers/deque.svh}}\par

Double-ended queue derived from vector. Supports push/pop at both ends, shuffle, and reverse.

\subsection*{\texttt{deque\#(T, P)}}
\noindent\small\textit{class}\quad \textbf{extends} \texttt{vector\#(T,P)}\par

Double-ended queue with push/pop at both front and back.

\begin{xltabular}{\linewidth}{>{\small}p{1.35in} >{\small}p{1.25in} >{\small}p{0.95in} >{\small}X}
\toprule
\textbf{Method} & \textbf{Return} & \textbf{Args} & \textbf{Description} \\
\midrule
\endhead
\bottomrule
\endfoot
\texttt{push\_front} & void & \texttt{T t} & Insert item at the front. \\
\texttt{push\_back} & void & \texttt{T t} & Insert item at the back. \\
\texttt{pop\_front} & \texttt{T} & --- & Remove and return front item. \\
\texttt{pop\_back} & \texttt{T} & --- & Remove and return back item. \\
\texttt{shuffle} & void & --- & Randomly reorder elements. \\
\texttt{reverse} & void & --- & Reverse the order of elements. \\
\texttt{clone} & \texttt{this\_t} & --- & Return a shallow copy of this deque. \\
\end{xltabular}

\section{circ\_buf --- Circular Buffer}
\noindent\small\textit{File: \texttt{src/containers/circ\_buf.svh}}\par

Fixed-capacity circular (ring) buffer backed by a static array of size S. Operates in reject mode (default) or overwrite mode. Parameterised on element type T, traits P, and compile-time size S.

\subsection*{\texttt{circ\_buf\#(T, P, S)}}
\noindent\small\textit{class}\quad \textbf{extends} \texttt{typed\_container\#(T,P)}\par

Fixed-size circular buffer. S elements, FIFO ordering.

\noindent\textbf{Constructor:} \texttt{new(---)} --- Calls clear() to initialise head, tail, and count.

\begin{xltabular}{\linewidth}{>{\small}p{1.35in} >{\small}p{1.25in} >{\small}p{0.95in} >{\small}X}
\toprule
\textbf{Method} & \textbf{Return} & \textbf{Args} & \textbf{Description} \\
\midrule
\endhead
\bottomrule
\endfoot
\texttt{size} & \texttt{size\_t} & --- & Number of items currently in the buffer. \\
\texttt{is\_empty} & \texttt{bit} & --- & True when count==0. \\
\texttt{is\_full} & \texttt{bit} & --- & True when count==S. \\
\texttt{clear} & void & --- & Reset head, tail, and count to zero. \\
\texttt{push} & void & \texttt{T t, circ\_buf\_mode\_e mode} & Insert t; in REJECT\_MODE prints error if full. In OVERWRITE\_MODE the tail item is overwritten. \\
\texttt{pop} & \texttt{T} & --- & Remove and return head item; prints error and returns empty if buffer is empty. \\
\end{xltabular}

\section{pair / triple / quadruple}
\noindent\small\textit{File: \texttt{src/containers/pair.svh}}\par

Lightweight tuple containers for binding two, three, or four values of potentially different types. All extend pair\#(T1,T2) and derive from void\_t (use void\_traits when placing in a container).

\subsection*{\texttt{pair\#(T1, T2)}}
\noindent\small\textit{class}\quad \textbf{extends} \texttt{void\_t}\par

Holds two values of types T1 and T2.

\noindent\textbf{Constructor:} \texttt{new(T1 a, T2 b)} --- Initialise both fields.

\begin{xltabular}{\linewidth}{>{\small}p{1.35in} >{\small}p{1.25in} >{\small}p{0.95in} >{\small}X}
\toprule
\textbf{Method} & \textbf{Return} & \textbf{Args} & \textbf{Description} \\
\midrule
\endhead
\bottomrule
\endfoot
\texttt{first} & \texttt{T1} & --- & Return the first element. \\
\texttt{second} & \texttt{T2} & --- & Return the second element. \\
\texttt{set\_first} & void & \texttt{T1 t} & Replace the first element. \\
\texttt{set\_second} & void & \texttt{T2 t} & Replace the second element. \\
\end{xltabular}

\subsection*{\texttt{triple\#(T1, T2, T3)}}
\noindent\small\textit{class}\quad \textbf{extends} \texttt{pair\#(T1,T2)}\par

Extends pair with a third element of type T3.

\noindent\textbf{Constructor:} \texttt{new(T1 a, T2 b, T3 c)} --- Initialise all three fields.

\begin{xltabular}{\linewidth}{>{\small}p{1.35in} >{\small}p{1.25in} >{\small}p{0.95in} >{\small}X}
\toprule
\textbf{Method} & \textbf{Return} & \textbf{Args} & \textbf{Description} \\
\midrule
\endhead
\bottomrule
\endfoot
\texttt{third} & \texttt{T3} & --- & Return the third element. \\
\texttt{set\_third} & void & \texttt{T3 t} & Replace the third element. \\
\end{xltabular}

\subsection*{\texttt{quadruple\#(T1, T2, T3, T4)}}
\noindent\small\textit{class}\quad \textbf{extends} \texttt{triple\#(T1,T2,T3)}\par

Extends triple with a fourth element of type T4.

\noindent\textbf{Constructor:} \texttt{new(T1 a, T2 b, T3 c, T4 d)} --- Initialise all four fields.

\begin{xltabular}{\linewidth}{>{\small}p{1.35in} >{\small}p{1.25in} >{\small}p{0.95in} >{\small}X}
\toprule
\textbf{Method} & \textbf{Return} & \textbf{Args} & \textbf{Description} \\
\midrule
\endhead
\bottomrule
\endfoot
\texttt{fourth} & \texttt{T4} & --- & Return the fourth element. \\
\texttt{set\_fourth} & void & \texttt{T4 t} & Replace the fourth element. \\
\end{xltabular}

\section{sorter}
\noindent\small\textit{File: \texttt{src/containers/sorter.svh}}\par

Generic quicksort implementation for SystemVerilog queues. Based on the algorithm in \textit{Algorithms in C} by Sedgewick. The comparator is provided by the traits parameter P.

\subsection*{\texttt{sorter\#(T, P)}}
\noindent\small\textit{class}\par

Quicksort for queues of type T using P::compare() as the comparator.

\begin{xltabular}{\linewidth}{>{\small}p{1.35in} >{\small}p{1.25in} >{\small}p{0.95in} >{\small}X}
\toprule
\textbf{Method} & \textbf{Return} & \textbf{Args} & \textbf{Description} \\
\midrule
\endhead
\bottomrule
\endfoot
\texttt{sort} & void & \texttt{ref T vec[\$]} & Public entry point: sort vec in place. Static. \\
\texttt{qsort} & void & \texttt{ref T vec[\$], int32\_t l, int32\_t r} & Recursive partition-and-sort step. Static. \\
\end{xltabular}

\chapter{Iterators}
\section{Iterator Interfaces}
\noindent\small\textit{File: \texttt{src/iterators/iterator\_intf.svh}}\par

The iterator interface hierarchy. \texttt{iterator\_intf\_base} defines set/get/skip/size/is\_empty. \texttt{fwd\_intf} adds forward traversal (first/next/is\_last/at\_end). \texttt{bkwd\_intf} adds backward traversal (last/prev/is\_first/at\_beginning). \texttt{random\_intf} adds random selection.

\subsection*{\texttt{iterator\_intf\_base\#(T, P)}}
\noindent\small\textit{interface class}\par

Root iterator interface. All iterators implement this.

\begin{xltabular}{\linewidth}{>{\small}p{1.35in} >{\small}p{1.25in} >{\small}p{0.95in} >{\small}X}
\toprule
\textbf{Method} & \textbf{Return} & \textbf{Args} & \textbf{Description} \\
\midrule
\endhead
\bottomrule
\endfoot
\texttt{set} & void & \texttt{T t} & Write value T at the current position (pure virtual). \\
\texttt{get} & \texttt{T} & --- & Read the value at the current position (pure virtual). \\
\texttt{skip} & \texttt{bit} & \texttt{signed\_index\_t d} & Move forward (d>0) or backward (d<0) by d steps (pure virtual). \\
\texttt{size} & \texttt{size\_t} & --- & Size of the bound container (pure virtual). \\
\texttt{is\_empty} & \texttt{bit} & --- & True if the iterator is unbound or the container is empty (pure virtual). \\
\end{xltabular}

\subsection*{\texttt{fwd\_intf\#(T, P)}}
\noindent\small\textit{interface class}\quad \textbf{extends} \texttt{iterator\_intf\_base\#(T,P)}\par

Forward-traversal interface.

\begin{xltabular}{\linewidth}{>{\small}p{1.35in} >{\small}p{1.25in} >{\small}p{0.95in} >{\small}X}
\toprule
\textbf{Method} & \textbf{Return} & \textbf{Args} & \textbf{Description} \\
\midrule
\endhead
\bottomrule
\endfoot
\texttt{first} & \texttt{bit} & --- & Reset to first element; returns 0 if empty (pure virtual). \\
\texttt{next} & \texttt{bit} & --- & Advance one step; may enter end state (pure virtual). \\
\texttt{is\_last} & \texttt{bit} & --- & True if pointing at the last element (pure virtual). \\
\texttt{at\_end} & \texttt{bit} & --- & True if past the last element (pure virtual). \\
\end{xltabular}

\subsection*{\texttt{bkwd\_intf\#(T, P)}}
\noindent\small\textit{interface class}\quad \textbf{extends} \texttt{iterator\_intf\_base\#(T,P)}\par

Backward-traversal interface.

\begin{xltabular}{\linewidth}{>{\small}p{1.35in} >{\small}p{1.25in} >{\small}p{0.95in} >{\small}X}
\toprule
\textbf{Method} & \textbf{Return} & \textbf{Args} & \textbf{Description} \\
\midrule
\endhead
\bottomrule
\endfoot
\texttt{last} & \texttt{bit} & --- & Reset to last element; returns 0 if empty (pure virtual). \\
\texttt{prev} & \texttt{bit} & --- & Move one step backward (pure virtual). \\
\texttt{is\_first} & \texttt{bit} & --- & True if pointing at the first element (pure virtual). \\
\texttt{at\_beginning} & \texttt{bit} & --- & True if before the first element (pure virtual). \\
\end{xltabular}

\subsection*{\texttt{random\_intf\#(T, P)}}
\noindent\small\textit{interface class}\quad \textbf{extends} \texttt{iterator\_intf\_base\#(T,P)}\par

Random-access iterator interface.

\begin{xltabular}{\linewidth}{>{\small}p{1.35in} >{\small}p{1.25in} >{\small}p{0.95in} >{\small}X}
\toprule
\textbf{Method} & \textbf{Return} & \textbf{Args} & \textbf{Description} \\
\midrule
\endhead
\bottomrule
\endfoot
\texttt{random} & \texttt{bit} & --- & Select a random valid element (pure virtual). \\
\texttt{set\_seed} & void & \texttt{int seed} & Set the RNG seed (pure virtual). \\
\texttt{set\_default\_seed} & void & --- & Restore the default RNG seed (pure virtual). \\
\end{xltabular}

\section{List Iterators}
\noindent\small\textit{File: \texttt{src/iterators/list\_iterators.svh}}\par

Iterators for vector and its subclasses (stack, queue, deque). \texttt{list\_iterator\_base} provides the common binding and set/get/skip/size/is\_empty logic. \texttt{list\_iterator} implements both forward and backward traversal. \texttt{list\_random\_iterator} implements random selection.

\subsection*{\texttt{list\_iterator\_base\#(T, P)}}
\noindent\small\textit{virtual class}\par

Abstract base for list iterators. Manages binding and common index operations.

\noindent\textbf{Constructor:} \texttt{new(list\_t list=null)} --- Optionally bind a vector.

\begin{xltabular}{\linewidth}{>{\small}p{1.35in} >{\small}p{1.25in} >{\small}p{0.95in} >{\small}X}
\toprule
\textbf{Method} & \textbf{Return} & \textbf{Args} & \textbf{Description} \\
\midrule
\endhead
\bottomrule
\endfoot
\texttt{bind\_list} & void & \texttt{list\_t list} & Bind iterator to a vector (or unbind if null). \\
\texttt{set} & void & \texttt{T t} & Write T at the current index position. \\
\texttt{get} & \texttt{T} & --- & Read value at the current index position. \\
\texttt{skip} & \texttt{bit} & \texttt{signed\_index\_t d} & Move index by d; clamps to valid range. \\
\texttt{size} & \texttt{size\_t} & --- & Return size of the bound vector. \\
\texttt{is\_empty} & \texttt{bit} & --- & True if unbound or vector is empty. \\
\end{xltabular}

\subsection*{\texttt{list\_iterator\#(T, P)}}
\noindent\small\textit{class}\quad \textbf{extends} \texttt{list\_iterator\_base\#(T,P)} \quad \textbf{implements} \texttt{fwd\_intf\#(T,P)}, \texttt{bkwd\_intf\#(T,P)}\par

Bidirectional iterator over any vector-derived container.

\noindent\textbf{Constructor:} \texttt{new(list\_t list=null)} --- Optionally bind a vector.

\begin{xltabular}{\linewidth}{>{\small}p{1.35in} >{\small}p{1.25in} >{\small}p{0.95in} >{\small}X}
\toprule
\textbf{Method} & \textbf{Return} & \textbf{Args} & \textbf{Description} \\
\midrule
\endhead
\bottomrule
\endfoot
\texttt{first} & \texttt{bit} & --- & Set index to 0; returns 0 if empty. \\
\texttt{next} & \texttt{bit} & --- & Increment index; returns 0 if past end. \\
\texttt{is\_last} & \texttt{bit} & --- & True if index is at the last valid element. \\
\texttt{at\_end} & \texttt{bit} & --- & True if index >= size. \\
\texttt{last} & \texttt{bit} & --- & Set index to last element; returns 0 if empty. \\
\texttt{prev} & \texttt{bit} & --- & Decrement index; returns 0 if already at start. \\
\texttt{is\_first} & \texttt{bit} & --- & True if index == 0. \\
\texttt{at\_beginning} & \texttt{bit} & --- & True if index < 0 (before first element). \\
\end{xltabular}

\subsection*{\texttt{list\_random\_iterator\#(T, P)}}
\noindent\small\textit{class}\quad \textbf{extends} \texttt{list\_iterator\_base\#(T,P)} \quad \textbf{implements} \texttt{random\_intf\#(T,P)}\par

Randomly selects an element from a bound vector.

\noindent\textbf{Constructor:} \texttt{new(list\_t list=null)} --- Optionally bind a vector; sets default RNG seed.

\begin{xltabular}{\linewidth}{>{\small}p{1.35in} >{\small}p{1.25in} >{\small}p{0.95in} >{\small}X}
\toprule
\textbf{Method} & \textbf{Return} & \textbf{Args} & \textbf{Description} \\
\midrule
\endhead
\bottomrule
\endfoot
\texttt{random} & \texttt{bit} & --- & Uniformly choose a random index in [0, size-1]. \\
\texttt{set\_seed} & void & \texttt{int seed} & Seed the RNG via \$urandom(seed). \\
\texttt{set\_default\_seed} & void & --- & Reset to the default seed (1). \\
\end{xltabular}

\section{Map Iterators}
\noindent\small\textit{File: \texttt{src/iterators/map\_iterators.svh}}\par

Iterators for the map container. \texttt{map\_iterator\_base} handles binding, index access, size, and is\_empty. \texttt{map\_iterator} provides bidirectional traversal. \texttt{map\_random\_iterator} provides random selection. Map keys are traversed in the natural order of the underlying SV associative array.

\subsection*{\texttt{map\_iterator\_base\#(KEY, T, P)}}
\noindent\small\textit{virtual class}\par

Abstract base for map iterators. Stores map binding and current key index.

\noindent\textbf{Constructor:} \texttt{new(map\_t map=null)} --- Optionally bind a map.

\begin{xltabular}{\linewidth}{>{\small}p{1.35in} >{\small}p{1.25in} >{\small}p{0.95in} >{\small}X}
\toprule
\textbf{Method} & \textbf{Return} & \textbf{Args} & \textbf{Description} \\
\midrule
\endhead
\bottomrule
\endfoot
\texttt{bind\_map} & void & \texttt{map\_t m} & Bind iterator to map m. \\
\texttt{get\_index} & \texttt{KEY} & --- & Return the current key. \\
\texttt{size} & \texttt{size\_t} & --- & Return number of entries in the bound map. \\
\texttt{is\_empty} & \texttt{bit} & --- & True if unbound or map is empty. \\
\end{xltabular}

\subsection*{\texttt{map\_iterator\#(KEY, T, P)}}
\noindent\small\textit{class}\quad \textbf{extends} \texttt{map\_iterator\_base\#(KEY,T,P)} \quad \textbf{implements} \texttt{fwd\_intf\#(T,P)}, \texttt{bkwd\_intf\#(T,P)}\par

Bidirectional iterator over a map. Tracks a four-state (INVALID/FIRST/VALID/LAST) traversal state.

\noindent\textbf{Constructor:} \texttt{new(map\_t map=null)} --- Optionally bind a map; set state to INVALID.

\begin{xltabular}{\linewidth}{>{\small}p{1.35in} >{\small}p{1.25in} >{\small}p{0.95in} >{\small}X}
\toprule
\textbf{Method} & \textbf{Return} & \textbf{Args} & \textbf{Description} \\
\midrule
\endhead
\bottomrule
\endfoot
\texttt{set} & void & \texttt{T t} & Update the value at the current key. \\
\texttt{get} & \texttt{T} & --- & Return the value at the current key. \\
\texttt{get\_index} & \texttt{KEY} & --- & Return the current key. \\
\texttt{first} & \texttt{bit} & --- & Move to first (smallest) key; set state FIRST. \\
\texttt{next} & \texttt{bit} & --- & Advance to next key; set state LAST if at end. \\
\texttt{is\_last} & \texttt{bit} & --- & True if current key equals the last key in the map. \\
\texttt{at\_end} & \texttt{bit} & --- & True when state==LAST. \\
\texttt{last} & \texttt{bit} & --- & Move to last key; set state LAST. \\
\texttt{prev} & \texttt{bit} & --- & Move to previous key; set state FIRST if at start. \\
\texttt{is\_first} & \texttt{bit} & --- & True if current key equals first key in the map. \\
\texttt{at\_beginning} & \texttt{bit} & --- & True when state==FIRST. \\
\texttt{skip} & \texttt{bit} & \texttt{signed\_index\_t d} & Skip d steps via repeated next()/prev() calls. \\
\end{xltabular}

\subsection*{\texttt{map\_random\_iterator\#(KEY, T, P)}}
\noindent\small\textit{class}\quad \textbf{extends} \texttt{map\_iterator\_base\#(KEY,T,P)} \quad \textbf{implements} \texttt{random\_intf\#(T,P)}\par

Randomly selects a map entry.

\noindent\textbf{Constructor:} \texttt{new(map\_t map=null)} --- Optionally bind a map; set default RNG seed.

\begin{xltabular}{\linewidth}{>{\small}p{1.35in} >{\small}p{1.25in} >{\small}p{0.95in} >{\small}X}
\toprule
\textbf{Method} & \textbf{Return} & \textbf{Args} & \textbf{Description} \\
\midrule
\endhead
\bottomrule
\endfoot
\texttt{set} & void & \texttt{T t} & Update value at current key. \\
\texttt{get} & \texttt{T} & --- & Return value at current key. \\
\texttt{random} & \texttt{bit} & --- & Select a uniformly random entry. \\
\texttt{set\_seed} & void & \texttt{int seed} & Seed the RNG via \$urandom(seed). \\
\texttt{set\_default\_seed} & void & --- & Reset to default seed (1). \\
\end{xltabular}

\section{Permutation Iterators}
\noindent\small\textit{File: \texttt{src/iterators/permute\_iterators.svh}}\par

Iterators that generate all permutations of a bound vector without modifying it. The permutation vector \texttt{pv} holds indices into the bound vector; \texttt{get\_nth(n)} returns the nth element for the current permutation. \texttt{permute\_iterator} traverses permutations sequentially; \texttt{permute\_random\_iterator} jumps to a randomly chosen permutation.

\subsection*{\texttt{permute\_iterator\_base\#(T, P)}}
\noindent\small\textit{class}\par

Shared permutation state and algorithm. Stores pv[], pix, and max\_pix.

\noindent\textbf{Constructor:} \texttt{new(vec\_t v=null)} --- Optionally bind a vector and initialise pv.

\begin{xltabular}{\linewidth}{>{\small}p{1.35in} >{\small}p{1.25in} >{\small}p{0.95in} >{\small}X}
\toprule
\textbf{Method} & \textbf{Return} & \textbf{Args} & \textbf{Description} \\
\midrule
\endhead
\bottomrule
\endfoot
\texttt{bind\_vector} & void & \texttt{vec\_t v} & Bind a vector; calls initialize(). \\
\texttt{size} & \texttt{size\_t} & --- & Size of the bound vector. \\
\texttt{is\_empty} & \texttt{bit} & --- & True if unbound or empty. \\
\texttt{get\_nth} & \texttt{T} & \texttt{index\_t n} & Return the nth element for the current permutation. \\
\texttt{get\_permutation\_index} & \texttt{longint} & --- & Return the current permutation index (0-based). \\
\texttt{get\_permutation\_vector} & \texttt{vector\#()} & --- & Return a copy of pv as a vector\#(index\_t, uint64\_traits). \\
\texttt{skip} & \texttt{bit} & \texttt{signed\_index\_t d} & Advance or retreat pix by d permutations. \\
\texttt{print} & void & --- & Display current pix and pv to stdout. \\
\end{xltabular}

\subsection*{\texttt{permute\_iterator\#(T, P)}}
\noindent\small\textit{class}\quad \textbf{extends} \texttt{permute\_iterator\_base\#(T,P)} \quad \textbf{implements} \texttt{fwd\_intf\#(T,P)}, \texttt{bkwd\_intf\#(T,P)}\par

Sequential bidirectional traversal of all permutations.

\noindent\textbf{Constructor:} \texttt{new(vec\_t v=null)} --- Optionally bind a vector.

\begin{xltabular}{\linewidth}{>{\small}p{1.35in} >{\small}p{1.25in} >{\small}p{0.95in} >{\small}X}
\toprule
\textbf{Method} & \textbf{Return} & \textbf{Args} & \textbf{Description} \\
\midrule
\endhead
\bottomrule
\endfoot
\texttt{first} & \texttt{bit} & --- & Reset to permutation 0. \\
\texttt{next} & \texttt{bit} & --- & Advance to next permutation; returns 0 if already at last. \\
\texttt{is\_last} & \texttt{bit} & --- & True if at the last permutation. \\
\texttt{at\_end} & \texttt{bit} & --- & True if pix >= max\textbackslash{}\_pix. \\
\texttt{last} & \texttt{bit} & --- & Jump to the last permutation. \\
\texttt{prev} & \texttt{bit} & --- & Return to the previous permutation. \\
\texttt{is\_first} & \texttt{bit} & --- & True if pix==0. \\
\texttt{at\_beginning} & \texttt{bit} & --- & True if pix<0. \\
\end{xltabular}

\subsection*{\texttt{permute\_random\_iterator\#(T, P)}}
\noindent\small\textit{class}\quad \textbf{extends} \texttt{permute\_iterator\_base\#(T,P)} \quad \textbf{implements} \texttt{random\_intf\#(T,P)}\par

Jump to a randomly chosen permutation.

\noindent\textbf{Constructor:} \texttt{new(vec\_t v=null)} --- Optionally bind a vector; set default RNG seed.

\begin{xltabular}{\linewidth}{>{\small}p{1.35in} >{\small}p{1.25in} >{\small}p{0.95in} >{\small}X}
\toprule
\textbf{Method} & \textbf{Return} & \textbf{Args} & \textbf{Description} \\
\midrule
\endhead
\bottomrule
\endfoot
\texttt{random} & \texttt{bit} & --- & Select a uniformly random permutation index. \\
\texttt{set\_seed} & void & \texttt{int seed} & Seed the RNG via \$urandom(seed). \\
\texttt{set\_default\_seed} & void & --- & Reset to default seed (1). \\
\end{xltabular}

\section{range}
\noindent\small\textit{File: \texttt{src/iterators/range.svh}}\par

A view over a contiguous sub-range [lb, ub] of an existing iterator. \texttt{range\_base} stores and validates the bounds. \texttt{range\#(T,P)} wraps a \texttt{fwd\_intf} iterator and optionally a \texttt{bkwd\_intf} if the underlying iterator supports it; backward access is silently unavailable otherwise.

\subsection*{\texttt{range\_base}}
\noindent\small\textit{virtual class}\par

Stores validated lower/upper index bounds. Swaps lb/ub if inverted.

\noindent\textbf{Constructor:} \texttt{new(size\_t size, index\_t lb, index\_t ub)} --- Clamp and optionally swap the bounds.

\begin{xltabular}{\linewidth}{>{\small}p{1.35in} >{\small}p{1.25in} >{\small}p{0.95in} >{\small}X}
\toprule
\textbf{Method} & \textbf{Return} & \textbf{Args} & \textbf{Description} \\
\midrule
\endhead
\bottomrule
\endfoot
\texttt{get\_lower\_bound} & \texttt{index\_t} & --- & Return lb. \\
\texttt{get\_upper\_bound} & \texttt{index\_t} & --- & Return ub. \\
\end{xltabular}

\subsection*{\texttt{range\#(T, P)}}
\noindent\small\textit{class}\quad \textbf{extends} \texttt{range\_base} \quad \textbf{implements} \texttt{fwd\_intf\#(T,P)}, \texttt{bkwd\_intf\#(T,P)}\par

Sub-range iterator. Binds to a fwd\_intf; optionally acquires backward capability via \$cast.

\noindent\textbf{Constructor:} \texttt{new(iterator\_intf\_base\#() it, index\_t lb, ub)} --- Cast it to fwd\_intf (fatal if fails); \$cast to bkwd\_intf (null if fails).

\begin{xltabular}{\linewidth}{>{\small}p{1.35in} >{\small}p{1.25in} >{\small}p{0.95in} >{\small}X}
\toprule
\textbf{Method} & \textbf{Return} & \textbf{Args} & \textbf{Description} \\
\midrule
\endhead
\bottomrule
\endfoot
\texttt{set} & void & \texttt{T t} & Set value at current iterator position. \\
\texttt{get} & \texttt{T} & --- & Get value at current iterator position. \\
\texttt{size} & \texttt{size\_t} & --- & Number of elements in the range: ub-lb+1. \\
\texttt{is\_empty} & \texttt{bit} & --- & True if underlying iterator is null or empty. \\
\texttt{first} & \texttt{bit} & --- & Move underlying iterator to position lb. \\
\texttt{next} & \texttt{bit} & --- & Advance one step within range bounds. \\
\texttt{is\_last} & \texttt{bit} & --- & True if current position >= ub. \\
\texttt{at\_end} & \texttt{bit} & --- & True if position > ub. \\
\texttt{last} & \texttt{bit} & --- & Move to position ub (requires bkwd\_intf). \\
\texttt{prev} & \texttt{bit} & --- & Move one step backward (requires bkwd\_intf). \\
\texttt{is\_first} & \texttt{bit} & --- & True if current position == lb. \\
\texttt{at\_beginning} & \texttt{bit} & --- & True if position < lb. \\
\texttt{skip} & \texttt{bit} & \texttt{signed\_index\_t d} & Delegate to underlying fwd\_iter.skip(d). \\
\end{xltabular}

\section{map\_range}
\noindent\small\textit{File: \texttt{src/iterators/map\_range.svh}}\par

Extends \texttt{range} for map iterators, adding \texttt{get\_index()} to retrieve the current map key from within a sub-range view.

\subsection*{\texttt{map\_range\#(KEY, T, P)}}
\noindent\small\textit{class}\quad \textbf{extends} \texttt{range\#(T,P)}\par

Range view over a map iterator slice, exposing the current key.

\noindent\textbf{Constructor:} \texttt{new(map\_iter\_t it, index\_t lb, ub)} --- Delegate to range; retain reference to the map iterator.

\begin{xltabular}{\linewidth}{>{\small}p{1.35in} >{\small}p{1.25in} >{\small}p{0.95in} >{\small}X}
\toprule
\textbf{Method} & \textbf{Return} & \textbf{Args} & \textbf{Description} \\
\midrule
\endhead
\bottomrule
\endfoot
\texttt{get\_index} & \texttt{KEY} & --- & Return the key at the current iterator position. \\
\end{xltabular}

\chapter{Algorithms}
\section{algo --- Core Algorithms}
\noindent\small\textit{File: \texttt{src/algorithms/algo.svh}}\par

Generic single-iterator algorithms. All functions are static and accept a \texttt{fwd\_intf} iterator plus, where needed, a \texttt{predicate} or \texttt{fcn} object. find() positions the iterator at the first matching element rather than returning it.

\subsection*{\texttt{algo\#(T, P)}}
\noindent\small\textit{class}\par

Collection of generic single-iterator algorithms.

\begin{xltabular}{\linewidth}{>{\small}p{1.35in} >{\small}p{1.25in} >{\small}p{0.95in} >{\small}X}
\toprule
\textbf{Method} & \textbf{Return} & \textbf{Args} & \textbf{Description} \\
\midrule
\endhead
\bottomrule
\endfoot
\texttt{count} & \texttt{uint32\_t} & \texttt{fwd\_intf\#() iter, predicate\#(T) p} & Count elements for which p.is\_true(). Static. \\
\texttt{all\_of} & \texttt{bit} & \texttt{fwd\_intf\#() iter, predicate\#(T) p} & True if predicate holds for every element. Static. \\
\texttt{none\_of} & \texttt{bit} & \texttt{fwd\_intf\#() iter, predicate\#(T) p} & True if predicate holds for no element. Static. \\
\texttt{any\_of} & \texttt{bit} & \texttt{fwd\_intf\#() iter, predicate\#(T) p} & True if predicate holds for at least one element. Static. \\
\texttt{min} & \texttt{T} & \texttt{fwd\_intf\#() iter} & Return minimum element using P::compare(). Static. \\
\texttt{max} & \texttt{T} & \texttt{fwd\_intf\#() iter} & Return maximum element using P::compare(). Static. \\
\texttt{find} & void & \texttt{fwd\_intf\#() iter, predicate\#(T) p} & Position iter at first element where p.is\_true(); no-op if not found. Static. \\
\texttt{for\_each} & void & \texttt{fwd\_intf\#() iter, fcn\#(T) fn} & Call fn.f(element) for every element. Static. \\
\end{xltabular}

\section{algo2 --- Dual-Iterator Algorithms}
\noindent\small\textit{File: \texttt{src/algorithms/algo2.svh}}\par

Algorithms that operate on two iterators simultaneously. \texttt{zip} traverses both iterators in lock-step, calling \texttt{fn.f(t1, t2)} at each step until either iterator reaches its end.

\subsection*{\texttt{algo2\#(T1, P1, T2, P2)}}
\noindent\small\textit{class}\par

Algorithms for paired traversal of two differently-typed iterators.

\begin{xltabular}{\linewidth}{>{\small}p{1.35in} >{\small}p{1.25in} >{\small}p{0.95in} >{\small}X}
\toprule
\textbf{Method} & \textbf{Return} & \textbf{Args} & \textbf{Description} \\
\midrule
\endhead
\bottomrule
\endfoot
\texttt{zip} & void & \texttt{fwd\_intf\#(T1,P1) iter1, fwd\_intf\#(T2,P2) iter2, fcn2\#(T1,T2) fn} & Traverse iter1 and iter2 in lock-step; call fn.f(t1,t2) at each step. Stops when either iterator ends. Static. \\
\end{xltabular}

\section{accum --- Accumulation Algorithms}
\noindent\small\textit{File: \texttt{src/algorithms/accum.svh}}\par

Algorithms that carry a mutable accumulator \texttt{A} across iterations. \texttt{accum} operates on one iterator; \texttt{accum2} operates on two in lock-step.

\subsection*{\texttt{accum\#(T, P, A)}}
\noindent\small\textit{class}\par

Single-iterator accumulation algorithm.

\begin{xltabular}{\linewidth}{>{\small}p{1.35in} >{\small}p{1.25in} >{\small}p{0.95in} >{\small}X}
\toprule
\textbf{Method} & \textbf{Return} & \textbf{Args} & \textbf{Description} \\
\midrule
\endhead
\bottomrule
\endfoot
\texttt{accumulate} & void & \texttt{fwd\_intf\#() iter, accum\_fcn\#(T,A) fn, ref A acc} & Call fn.f(element, acc) for every element; acc is updated in place by fn. Static. \\
\end{xltabular}

\subsection*{\texttt{accum2\#(T1, P1, T2, P2, A)}}
\noindent\small\textit{class}\par

Dual-iterator accumulation algorithm.

\begin{xltabular}{\linewidth}{>{\small}p{1.35in} >{\small}p{1.25in} >{\small}p{0.95in} >{\small}X}
\toprule
\textbf{Method} & \textbf{Return} & \textbf{Args} & \textbf{Description} \\
\midrule
\endhead
\bottomrule
\endfoot
\texttt{accumulate} & void & \texttt{fwd\_intf\#(T1,P1) iter1, fwd\_intf\#(T2,P2) iter2, accum\_fcn2\#() fn, ref A acc} & Traverse both iterators in lock-step; call fn.f(t1, t2, acc) at each step. Static. \\
\end{xltabular}

\section{fcn --- Function Objects}
\noindent\small\textit{File: \texttt{src/algorithms/fcn.svh}}\par

Virtual function objects (SVX's substitute for C++ functors and lambdas). Users derive concrete classes and override \texttt{f()}. The hierarchy provides one-argument, two-argument, and accumulating variants.

\subsection*{\texttt{fcn\_base}}
\noindent\small\textit{virtual class}\par

Empty base class for all function objects.


\subsection*{\texttt{fcn\#(T)}}
\noindent\small\textit{virtual class}\quad \textbf{extends} \texttt{fcn\_base}\par

Abstract function object taking one argument of type T.

\begin{xltabular}{\linewidth}{>{\small}p{1.35in} >{\small}p{1.25in} >{\small}p{0.95in} >{\small}X}
\toprule
\textbf{Method} & \textbf{Return} & \textbf{Args} & \textbf{Description} \\
\midrule
\endhead
\bottomrule
\endfoot
\texttt{f} & void & \texttt{T t} & Apply the function to t (pure virtual). \\
\end{xltabular}

\subsection*{\texttt{fcn2\#(T1, T2)}}
\noindent\small\textit{virtual class}\quad \textbf{extends} \texttt{fcn\_base}\par

Abstract function object taking two arguments of types T1 and T2.

\begin{xltabular}{\linewidth}{>{\small}p{1.35in} >{\small}p{1.25in} >{\small}p{0.95in} >{\small}X}
\toprule
\textbf{Method} & \textbf{Return} & \textbf{Args} & \textbf{Description} \\
\midrule
\endhead
\bottomrule
\endfoot
\texttt{f} & void & \texttt{T1 t1, T2 t2} & Apply the function to (t1, t2) (pure virtual). \\
\end{xltabular}

\subsection*{\texttt{accum\_fcn\#(T, A)}}
\noindent\small\textit{virtual class}\quad \textbf{extends} \texttt{fcn\_base}\par

Abstract function object taking one value and a mutable accumulator.

\begin{xltabular}{\linewidth}{>{\small}p{1.35in} >{\small}p{1.25in} >{\small}p{0.95in} >{\small}X}
\toprule
\textbf{Method} & \textbf{Return} & \textbf{Args} & \textbf{Description} \\
\midrule
\endhead
\bottomrule
\endfoot
\texttt{f} & void & \texttt{T t, ref A a} & Apply the function; update accumulator a in place (pure virtual). \\
\end{xltabular}

\subsection*{\texttt{accum\_fcn2\#(T1, T2, A)}}
\noindent\small\textit{virtual class}\quad \textbf{extends} \texttt{fcn\_base}\par

Abstract function object taking two values and a mutable accumulator.

\begin{xltabular}{\linewidth}{>{\small}p{1.35in} >{\small}p{1.25in} >{\small}p{0.95in} >{\small}X}
\toprule
\textbf{Method} & \textbf{Return} & \textbf{Args} & \textbf{Description} \\
\midrule
\endhead
\bottomrule
\endfoot
\texttt{f} & void & \texttt{T1 t1, T2 t2, ref A a} & Apply the function to (t1, t2); update a in place (pure virtual). \\
\end{xltabular}

\section{predicate --- Predicate Classes}
\noindent\small\textit{File: \texttt{src/algorithms/predicate.svh}}\par

Virtual predicate hierarchy. A predicate is a Boolean function over a single element T. Concrete predicate classes are composed with \texttt{and\_pred}, \texttt{or\_pred}, and \texttt{not\_pred}.

\subsection*{\texttt{predicate\_base}}
\noindent\small\textit{virtual class}\par

Empty root for all predicate classes.


\subsection*{\texttt{predicate\#(T)}}
\noindent\small\textit{virtual class}\quad \textbf{extends} \texttt{predicate\_base}\par

Abstract predicate: is\_true() is pure virtual; is\_false() delegates.

\begin{xltabular}{\linewidth}{>{\small}p{1.35in} >{\small}p{1.25in} >{\small}p{0.95in} >{\small}X}
\toprule
\textbf{Method} & \textbf{Return} & \textbf{Args} & \textbf{Description} \\
\midrule
\endhead
\bottomrule
\endfoot
\texttt{is\_true} & \texttt{bit} & \texttt{T t} & Return true if t satisfies the predicate (pure virtual). \\
\texttt{is\_false} & \texttt{bit} & \texttt{T t} & Return !is\_true(t). \\
\end{xltabular}

\subsection*{\texttt{and\_pred\#(T)}}
\noindent\small\textit{class}\quad \textbf{extends} \texttt{predicate\#(T)}\par

True if both sub-predicates a and b are true.

\noindent\textbf{Constructor:} \texttt{new(predicate\#(T) \_a, \_b)} --- Store references to a and b.

\begin{xltabular}{\linewidth}{>{\small}p{1.35in} >{\small}p{1.25in} >{\small}p{0.95in} >{\small}X}
\toprule
\textbf{Method} & \textbf{Return} & \textbf{Args} & \textbf{Description} \\
\midrule
\endhead
\bottomrule
\endfoot
\texttt{is\_true} & \texttt{bit} & \texttt{T t} & Return a.is\_true(t) \&\& b.is\_true(t). \\
\end{xltabular}

\subsection*{\texttt{or\_pred\#(T)}}
\noindent\small\textit{class}\quad \textbf{extends} \texttt{predicate\#(T)}\par

True if at least one of the two sub-predicates is true.

\noindent\textbf{Constructor:} \texttt{new(predicate\#(T) \_a, \_b)} --- Store references to a and b.

\begin{xltabular}{\linewidth}{>{\small}p{1.35in} >{\small}p{1.25in} >{\small}p{0.95in} >{\small}X}
\toprule
\textbf{Method} & \textbf{Return} & \textbf{Args} & \textbf{Description} \\
\midrule
\endhead
\bottomrule
\endfoot
\texttt{is\_true} & \texttt{bit} & \texttt{T t} & Return a.is\_true(t) || b.is\_true(t). \\
\end{xltabular}

\subsection*{\texttt{not\_pred\#(T)}}
\noindent\small\textit{class}\quad \textbf{extends} \texttt{predicate\#(T)}\par

True when the wrapped predicate is false.

\noindent\textbf{Constructor:} \texttt{new(predicate\#(T) \_p)} --- Store reference to p.

\begin{xltabular}{\linewidth}{>{\small}p{1.35in} >{\small}p{1.25in} >{\small}p{0.95in} >{\small}X}
\toprule
\textbf{Method} & \textbf{Return} & \textbf{Args} & \textbf{Description} \\
\midrule
\endhead
\bottomrule
\endfoot
\texttt{is\_true} & \texttt{bit} & \texttt{T t} & Return p.is\_false(t). \\
\end{xltabular}

\subsection*{\texttt{always\_true\#(T)}}
\noindent\small\textit{class}\quad \textbf{extends} \texttt{predicate\#(T)}\par

Constant predicate that always evaluates to 1.

\begin{xltabular}{\linewidth}{>{\small}p{1.35in} >{\small}p{1.25in} >{\small}p{0.95in} >{\small}X}
\toprule
\textbf{Method} & \textbf{Return} & \textbf{Args} & \textbf{Description} \\
\midrule
\endhead
\bottomrule
\endfoot
\texttt{is\_true} & \texttt{bit} & \texttt{T t} & Always returns 1. \\
\end{xltabular}

\subsection*{\texttt{always\_false\#(T)}}
\noindent\small\textit{class}\quad \textbf{extends} \texttt{predicate\#(T)}\par

Constant predicate that always evaluates to 0.

\begin{xltabular}{\linewidth}{>{\small}p{1.35in} >{\small}p{1.25in} >{\small}p{0.95in} >{\small}X}
\toprule
\textbf{Method} & \textbf{Return} & \textbf{Args} & \textbf{Description} \\
\midrule
\endhead
\bottomrule
\endfoot
\texttt{is\_true} & \texttt{bit} & \texttt{T t} & Always returns 0. \\
\end{xltabular}

\chapter{Linked Structures}
\section{node}
\noindent\small\textit{File: \texttt{src/linked/node.svh}}\par

A named, markable container node extending \texttt{object}. The name field supports hierarchical naming in the tree. The mark/unmark mechanism supports traversal algorithms that need to track visited nodes.

\subsection*{\texttt{node}}
\noindent\small\textit{class}\quad \textbf{extends} \texttt{object}\par

Named, markable base node for tree structures.

\noindent\textbf{Constructor:} \texttt{new(string nm="")} --- Initialise name to nm.

\begin{xltabular}{\linewidth}{>{\small}p{1.35in} >{\small}p{1.25in} >{\small}p{0.95in} >{\small}X}
\toprule
\textbf{Method} & \textbf{Return} & \textbf{Args} & \textbf{Description} \\
\midrule
\endhead
\bottomrule
\endfoot
\texttt{set\_name} & void & \texttt{string nm} & Set the node name. \\
\texttt{get\_name} & \texttt{string} & --- & Return the node name. \\
\texttt{mark} & void & --- & Set the mark flag to 1. \\
\texttt{unmark} & void & --- & Clear the mark flag to 0. \\
\texttt{get\_mark} & \texttt{bit} & --- & Return the current mark value. \\
\texttt{is\_marked} & --- & --- & Return the mark flag (equivalent to get\_mark()). \\
\end{xltabular}

\section{tree}
\noindent\small\textit{File: \texttt{src/linked/tree.svh}}\par

Classic tree data structure. Each node has a name, a reference to its parent, and a map of named children. Hierarchical path names use the dot (.) separator. The \texttt{find()} method parses a dot-separated path using the lexer and locates the named node.

\subsection*{\texttt{tree}}
\noindent\small\textit{class}\quad \textbf{extends} \texttt{node}\par

N-ary tree node with named children, full-path support, and hierarchical search.

\noindent\textbf{Constructor:} \texttt{new(string nm, tree p)} --- Set name, link to parent p; register as child of p; build full\_name.

\begin{xltabular}{\linewidth}{>{\small}p{1.35in} >{\small}p{1.25in} >{\small}p{0.95in} >{\small}X}
\toprule
\textbf{Method} & \textbf{Return} & \textbf{Args} & \textbf{Description} \\
\midrule
\endhead
\bottomrule
\endfoot
\texttt{insert} & void & \texttt{tree t} & Add t as a named child of this node; error if name duplicate. \\
\texttt{get\_parent} & \texttt{tree} & --- & Return the parent node handle. \\
\texttt{get\_full\_name} & \texttt{string} & --- & Return the fully-qualified dotted path name. \\
\texttt{first\_child} & \texttt{bit} & \texttt{ref string nm} & Set nm to the name of the first child; returns 0 if none. \\
\texttt{next\_child} & \texttt{bit} & \texttt{ref string nm} & Advance nm to the next sibling name; returns 0 if none. \\
\texttt{num\_children} & \texttt{size\_t} & --- & Return the number of direct children. \\
\texttt{size} & \texttt{size\_t} & --- & Return the number of direct children (delegates to get\_children()). \\
\texttt{get\_children} & \texttt{deque\#()} & --- & Return a deque of all direct child handles. \\
\texttt{get\_child} & \texttt{tree} & \texttt{string nm} & Return the named child; null if not found. \\
\texttt{find} & \texttt{tree} & \texttt{string path} & Locate a node by dot-separated path. Returns null on parse error or not found. \\
\texttt{mark\_all} & void & --- & Recursively mark all nodes in this subtree. \\
\texttt{unmark\_all} & void & --- & Recursively unmark all nodes in this subtree. \\
\end{xltabular}

\section{tree\_iterator}
\noindent\small\textit{File: \texttt{src/linked/tree\_iterator.svh}}\par

Pre-order depth-first iterator over a tree. A stack drives the traversal: \texttt{first()} seeds the stack with the root; \texttt{next()} pops the current node and pushes its children. \texttt{size()} counts all reachable nodes using a separate DFS with marks.

\subsection*{\texttt{tree\_iterator\#(T, P)}}
\noindent\small\textit{class}\quad \textbf{implements} \texttt{fwd\_intf\#(T,P)}\par

Forward-only pre-order DFS iterator over a tree.

\noindent\textbf{Constructor:} \texttt{new(tree t=null)} --- Bind to tree t; initialise internal stack.

\begin{xltabular}{\linewidth}{>{\small}p{1.35in} >{\small}p{1.25in} >{\small}p{0.95in} >{\small}X}
\toprule
\textbf{Method} & \textbf{Return} & \textbf{Args} & \textbf{Description} \\
\midrule
\endhead
\bottomrule
\endfoot
\texttt{size} & \texttt{size\_t} & --- & Count all nodes with a separate DFS using marks. Side effect: alters node marks. \\
\texttt{is\_empty} & \texttt{bit} & --- & True if m\_tree is null. \\
\texttt{set} & void & \texttt{tree t} & No-op (required by fwd\_intf; not applicable to tree traversal). \\
\texttt{get} & \texttt{tree} & --- & Return the tree node at the top of the stack. \\
\texttt{first} & \texttt{bit} & --- & Reset stack; push root; unmark all nodes. \\
\texttt{next} & \texttt{bit} & --- & Pop current node; push all its children onto the stack. \\
\texttt{is\_last} & \texttt{bit} & --- & True when exactly one node remains on the stack. \\
\texttt{at\_end} & \texttt{bit} & --- & True when the stack is empty. \\
\texttt{skip} & \texttt{bit} & \texttt{signed\_index\_t d} & Skip d nodes forward (negative d not supported). \\
\end{xltabular}

\chapter{Set}
\section{set}
\noindent\small\textit{File: \texttt{src/set/set\_impl.svh}}\par

Mathematical set backed by a \texttt{map\#(T, uint32\_t, uint32\_traits)}. Because map keys are maintained in sorted order by the SV associative array, set operations such as intersection use a sorted merge-walk (O(n+m)) rather than membership tests. All set operations return a new set and do not mutate their operands.

\subsection*{\texttt{set\#(T, P)}}
\noindent\small\textit{class}\quad \textbf{extends} \texttt{typed\_container\#(T,P)}\par

Mathematical set with insert, membership test, and set-algebra operations.

\noindent\textbf{Constructor:} \texttt{new(---)} --- Allocate the internal map.

\begin{xltabular}{\linewidth}{>{\small}p{1.35in} >{\small}p{1.25in} >{\small}p{0.95in} >{\small}X}
\toprule
\textbf{Method} & \textbf{Return} & \textbf{Args} & \textbf{Description} \\
\midrule
\endhead
\bottomrule
\endfoot
\texttt{size} & \texttt{size\_t} & --- & Return the number of members. \\
\texttt{is\_empty} & \texttt{bit} & --- & True when the set has no members. \\
\texttt{clear} & void & --- & Remove all members. \\
\texttt{insert} & void & \texttt{T item} & Add item to the set (duplicates ignored via map semantics). \\
\texttt{contains} & \texttt{bit} & \texttt{T item} & True if item is a member. \\
\texttt{intersection} & \texttt{set\_t} & \texttt{set\_t s} & Return a new set of elements in both this and s. Uses sorted merge-walk. \\
\texttt{set\_union} & \texttt{set\_t} & \texttt{set\_t s} & Return a new set of all elements in this or s (or both). \\
\texttt{difference} & \texttt{set\_t} & \texttt{set\_t s} & Return a new set of elements in this but not in s. \\
\texttt{to\_vector} & \texttt{vector\#()} & --- & Return a vector containing all set members in key order. \\
\texttt{from\_vector} & void & \texttt{vector\#() vec} & Insert all elements of vec into the set (skipping duplicates). \\
\end{xltabular}

\end{document}
